\documentclass[10pt,a4paper]{amsart}
\usepackage{amsmath,amssymb,amsthm,graphicx}
\usepackage{verbatim}
\newenvironment{reference}{\comment}{\endcomment}
%\newenvironment{reference}{}{}
\newenvironment{slogan}{\comment}{\endcomment}
\newenvironment{history}{\comment}{\endcomment}

% For commutative diagrams we use Xy-pic
\usepackage[all]{xy}

% We use 2cell for 2-commutative diagrams.
\xyoption{2cell}
% Theorem environments.
%
\theoremstyle{plain}
\newtheorem{theorem}[subsection]{Theorem}
\newtheorem{proposition}[subsection]{Proposition}
\newtheorem{lemma}[subsection]{Lemma}

\theoremstyle{definition}
\newtheorem{definition}[subsection]{Definition}
\newtheorem{example}[subsection]{Example}
\newtheorem{exercise}[subsection]{Exercise}
\newtheorem{situation}[subsection]{Situation}

\theoremstyle{remark}
\newtheorem{remark}[subsection]{Remark}
\newtheorem{remarks}[subsection]{Remarks}

\numberwithin{equation}{subsection}

% Macros
%
\def\lim{\mathop{\mathrm{lim}}\nolimits}
\def\colim{\mathop{\mathrm{colim}}\nolimits}
\def\Spec{\mathop{\mathrm{Spec}}}
\def\Hom{\mathop{\mathrm{Hom}}\nolimits}
\def\Ext{\mathop{\mathrm{Ext}}\nolimits}
\def\SheafHom{\mathop{\mathcal{H}\!\mathit{om}}\nolimits}
\def\SheafExt{\mathop{\mathcal{E}\!\mathit{xt}}\nolimits}
\def\Sch{\mathit{Sch}}
\def\Mor{\mathop{\mathrm{Mor}}\nolimits}
\def\Ob{\mathop{\mathrm{Ob}}\nolimits}
\def\Sh{\mathop{\mathit{Sh}}\nolimits}
\def\NL{\mathop{N\!L}\nolimits}
\def\CH{\mathop{\mathrm{CH}}\nolimits}
\def\proetale{{pro\text{-}\acute{e}tale}}
\def\etale{{\acute{e}tale}}
\def\QCoh{\mathit{QCoh}}
\def\Ker{\mathop{\mathrm{Ker}}}
\def\Im{\mathop{\mathrm{Im}}}
\def\Coker{\mathop{\mathrm{Coker}}}
\def\Coim{\mathop{\mathrm{Coim}}}

% Boxtimes
%
\DeclareMathSymbol{\boxtimes}{\mathbin}{AMSa}{"02}

%
% Macros for moduli stacks/spaces
%
\def\QCohstack{\mathcal{QC}\!\mathit{oh}}
\def\Cohstack{\mathcal{C}\!\mathit{oh}}
\def\Spacesstack{\mathcal{S}\!\mathit{paces}}
\def\Quotfunctor{\mathrm{Quot}}
\def\Hilbfunctor{\mathrm{Hilb}}
\def\Curvesstack{\mathcal{C}\!\mathit{urves}}
\def\Polarizedstack{\mathcal{P}\!\mathit{olarized}}
\def\Complexesstack{\mathcal{C}\!\mathit{omplexes}}
% \Pic is the operator that assigns to X its picard group, usage \Pic(X)
% \Picardstack_{X/B} denotes the Picard stack of X over B
% \Picardfunctor_{X/B} denotes the Picard functor of X over B
\def\Pic{\mathop{\mathrm{Pic}}\nolimits}
\def\Picardstack{\mathcal{P}\!\mathit{ic}}
\def\Picardfunctor{\mathrm{Pic}}
\def\Deformationcategory{\mathcal{D}\!\mathit{ef}}
\usepackage{cleveref}
\crefname{originalreference}{original reference}{original references}
\Crefname{originalreference}{Original reference}{Original references}
\makeatletter
\newlabel{lemma-chow-ring-projective-bundle}{{02TX}{source}{}{sourcepointer-1}{}}
\newlabel{lemma-chow-ring-projective-bundle@cref}{{[originalreference][0][]02TX}{[1][0][]source}{}{}{}}
\makeatother
\crefname{originalreference}{original reference}{original references}
\Crefname{originalreference}{Original reference}{Original references}
\makeatletter
\newlabel{situation-setup}{{02QL}{source}{}{sourcepointer-2}{}}
\newlabel{situation-setup@cref}{{[originalreference][0][]02QL}{[1][0][]source}{}{}{}}
\makeatother
\title{Splitting principle: a projective flat cover trivialises a finite collection of locally free sheaves and is injective on Chow groups}
\author{The Stacks Project Authors}
\date{stacks-project:spaces-chow.tex:local-snapshot-20261012; GNU FDL 1.2 or later}
\begin{document}
\maketitle
\noindent\textit{Editorial note.} All delivered mathematical text is the author's own, copied from the pinned Stacks Project chapter file spaces-chow.tex; no statement or proof step was rewritten, reordered or removed.\par\medskip
\noindent\textit{Editorial note.} The counted proof is the author's complete argument: an induction on the total rank, passing to the projective bundle of a sheaf of larger rank and replacing it by the locally free kernel of the canonical surjection\par\medskip
\noindent\textit{Editorial note.} The label lemma-splitting-principle names the counted statement, and the proof environment immediately following it in the source is delivered verbatim as its proof.\par\medskip
\noindent\textit{Editorial note.} The counted statement is carried by the author's own lemma environment in the chapter source.\par\medskip
\noindent\textit{Allowed standard prerequisites (editorial).}\begin{itemize}
\item The projective bundle formula and the associated canonical surjection onto the tautological quotient with its finite locally free kernel of rank one less
\item Injectivity of pullback along the projection of a projective bundle, for every base change
\item Filtrations of locally free sheaves with invertible successive quotients, and flat projective morphisms of a given relative dimension
\end{itemize}\medskip
\ifdefined\hypertarget\hypertarget{sourcepointer-1}{}\fi
\noindent\textit{Source reference.} Original reference 02TX: sources/stacks-spaces-chow/pinned-original-source.tex, source lines 4381--4381. Prerequisite pointer invoked at proof line(s) 5031; it names a standard fact of the same chapter that the counted argument uses, and every step of the counted proof built on it is retained verbatim in the delivered body.\par\medskip
\ifdefined\hypertarget\hypertarget{sourcepointer-2}{}\fi
\noindent\textit{Source reference.} Original reference 02QL: sources/stacks-spaces-chow/pinned-original-source.tex, source lines 264--264. Prerequisite pointer invoked at proof line(s) 5005; it names a standard fact of the same chapter that the counted argument uses, and every step of the counted proof built on it is retained verbatim in the delivered body.\par\medskip
\begin{lemma}
\label{lemma-splitting-principle}
In Situation \ref{situation-setup} let $X/B$ be good.
Let $\mathcal{E}_i$ be a finite collection of locally free
$\mathcal{O}_X$-modules of rank $r_i$. There exists a projective
flat morphism $\pi : P \to X$ of relative dimension $d$ such that
\begin{enumerate}
\item for any morphism $f : Y \to X$ of good algebraic spaces
over $B$ the map
$\pi_Y^* : \CH_*(Y) \to \CH_{* + d}(Y \times_X P)$ is injective, and
\item each $\pi^*\mathcal{E}_i$ has a filtration
whose successive quotients $\mathcal{L}_{i, 1}, \ldots, \mathcal{L}_{i, r_i}$
are invertible ${\mathcal O}_P$-modules.
\end{enumerate}
\end{lemma}

\begin{proof}
We prove this by induction on the integer $r = \sum r_i$.
If $r = 0$ we can take $\pi = \text{id}_X$.
If $r_i = 1$ for all $i$, then we can also take $\pi = \text{id}_X$.
Assume that $r_{i_0} > 1$ for some $i_0$.
Let $(\pi : P \to X, \mathcal{O}_P(1))$
be the projective bundle associated to $\mathcal{E}_{i_0}$.
The canonical map $\pi^*\mathcal{E}_{i_0} \to \mathcal{O}_P(1)$
is surjective and hence its kernel $\mathcal{E}'_{i_0}$
is finite locally free of rank $r_{i_0} - 1$.
Observe that $\pi_Y^*$ is injective for any
morphism $f : Y \to X$ of good algebraic spaces
over $B$, see Lemma \ref{lemma-chow-ring-projective-bundle}.
Thus it suffices to prove the lemma for
$P$ and the locally free sheaves $\pi^*\mathcal{E}_i$.
However, because we have the subbundle
$\mathcal{E}_{i_0} \subset \pi^*\mathcal{E}_{i_0}$
with invertible quotient, it now suffices to prove
the lemma for the collection
$\{\mathcal{E}_i\}_{i \not = i_0} \cup \{\mathcal{E}'_{i_0}\}$.
This decreases $r$ by $1$ and we win by induction hypothesis.
\end{proof}

\end{document}
